\documentclass[11pt]{article}
\usepackage[margin=2.5cm]{geometry}
\usepackage{amsmath}
\usepackage{parskip}

\newcommand{\avg}[1]{\langle #1 \rangle}

\title{\textbf{Processing plan: cross-wind PIV turbulent statistics}\\[4pt]
       \large Delaware wind-wave-current tank, JFM 2023 case}
\author{}
\date{\today}

\begin{document}
\maketitle
\thispagestyle{empty}

\textbf{Goal.} Compute horizontally averaged turbulent statistics
$\avg{w'w'}(z,t)$, $\avg{v'v'}(z,t)$, $\avg{v'w'}(z,t)$
from the cross-wind (transverse) PIV data (Ramp 2, three repeats),
for direct comparison with Greg's Oceananigans LES.

\section*{The wave removal problem}

The raw PIV velocity field contains both wave orbital motion and turbulence.
Before computing turbulent statistics I need to isolate the turbulent part:
\[
  w_\text{meas}(y,z,t) = w_\text{wave}(z,t) + w'(y,z,t).
\]
Since the PIV plane is cross-wind $(y,z)$ and waves propagate along-wind,
the wave orbital velocity is \emph{uniform across} $y$.
The orbital $w$ decays exponentially with depth as $e^{kz}$; for
capillary ripples ($k \approx 209$~m$^{-1}$, $\lambda \approx 3$~cm)
the e-folding depth is only $\sim$5~mm, so wave contamination is
confined to the top $\sim$1~cm.

\section*{How I compute $w_\text{wave}(z,t)$}

I use $\eta(x,t)$ extracted from the longitudinal experiment
(\texttt{ExpLCL\_2\_01}) to get the local wave properties at each time step.

\begin{enumerate}
  \item Apply the Hilbert transform to $\eta(x,t)$ along $x$ at each $t$
        to get the analytic signal $\hat\eta = a\,e^{i\varphi}$.
  \item Smooth the unwrapped phase $\varphi$ and differentiate to get the
        local wavenumber $k(x,t) = \partial\varphi/\partial x$.
  \item Convert to frequency via the gravity-capillary dispersion relation
        (deep water, $kH \approx 148 \gg 1$):
        \[
          \omega = \sqrt{k(g + \gamma k^2)}, \qquad
          \gamma = 7.2\times10^{-5}~\text{m}^3\,\text{s}^{-2}.
        \]
  \item The vertical orbital velocity at depth is then
        \[
          w_\text{wave}(z,t) = a\,\omega\,e^{kz}\sin\varphi
          \Big|_{x\,=\,12~\text{m}}.
        \]
\end{enumerate}

I then subtract $w_\text{wave}(z,t)$ from each column of the PIV frame
(it is uniform in $y$) to obtain $w'(y,z,t)$.

\section*{Alternative: $y$-mean subtraction}

Because $w_\text{wave}$ is uniform in $y$, the horizontal mean
$\avg{w_\text{meas}}_y$ directly captures the wave signal, so
\[
  w'(y,z,t) = w_\text{meas}(y,z,t) - \avg{w_\text{meas}}_y(z,t)
\]
is a simpler alternative that needs no $\eta$ data.
I will run both methods and compare: they should agree below
$|z| \approx 1$~cm where the orbital velocity has decayed to negligible levels.

\section*{Turbulent statistics}

After wave removal I compute the $y$-averaged second moments at each
$(z,t)$:
\[
  \avg{w'w'},\quad \avg{v'v'},\quad \avg{v'w'},\quad \avg{w'^3}.
\]
The result is a depth--time Hovm\"{o}ller diagram comparable to
Greg's LES output and to JFM Fig.~8.

\section*{Open questions for the group}

\begin{enumerate}

  \item \textbf{Time alignment.}
  The $\eta$ data and the transverse PIV start at slightly different
  absolute times and may drift.
  How careful do I need to be about aligning them in time ---
  is a simple linear interpolation onto the PIV time vector enough,
  or does a small timing offset matter for the wave subtraction?

  \item \textbf{Hilbert transform and ``local'' wave properties.}
  The Hilbert transform gives an instantaneous wavenumber $k(x,t)$
  that varies in space and time rather than using a single
  peak wavenumber from a spectrum.
  Is this the right thing to use here, or is the wave field
  narrow-band enough that a single $k$ from the dispersion peak
  would be just as good?

  \item \textbf{Hilbert vs.\ $y$-mean subtraction.}
  The simpler method (subtract $\avg{w}_y$) exploits the fact that
  the wave is uniform across $y$ and needs no $\eta$ data at all.
  Is there a reason to prefer the Hilbert approach, or is
  $y$-mean subtraction sufficient given that orbital velocities
  decay to $<5\%$ below 1~cm depth?

  \item \textbf{Why these statistics, and what should I expect?}
  $\avg{w'w'}$ should increase sharply at $t\approx16$--18~s
  (Langmuir instability onset) and deepen over time.
  But the PIV window starts at $t=38$~s --- so the transition
  itself is not captured.
  Is Greg primarily interested in the \emph{developed} turbulence
  profiles (38--98~s), and if so, what functional form does he
  expect for $\avg{w'w'}(z)$?

  \item \textbf{What else to plot?}
  Beyond the Hovm\"{o}ller diagram, would it be useful to show
  (a) vertical profiles $\avg{w'w'}(z)$ at fixed times,
  (b) time series of $\avg{w'w'}$ at a fixed depth,
  or (c) the ratio $\avg{w'w'}/\avg{v'v'}$ as an anisotropy
  diagnostic?
  What format does Greg use in the LES output so the axes match?

\end{enumerate}

\end{document}
